
\documentclass{article}
\usepackage{colortbl}
\usepackage{makecell}
\usepackage{multirow}
\usepackage{supertabular}

\begin{document}

\newcounter{utterance}

\centering \large Interaction Transcript for game `cladder', experiment `full\_v1.5\_default', episode 319 with llama{-}3.1{-}8b{-}instruct.
\vspace{24pt}

{ \footnotesize  \setcounter{utterance}{1}
\setlength{\tabcolsep}{0pt}
\begin{supertabular}{c@{$\;$}|p{.15\linewidth}@{}p{.15\linewidth}p{.15\linewidth}p{.15\linewidth}p{.15\linewidth}p{.15\linewidth}}
    \# & \multicolumn{2}{c}{Player} && \multicolumn{2}{c}{Game Master} \\
    \hline

    \theutterance \stepcounter{utterance}  
    & & & \multicolumn{4}{p{0.6\linewidth}}{
        \cellcolor[rgb]{0.9,0.9,0.9}{
            \makecell[{{p{\linewidth}}}]{
                \texttt{\tiny{[P1$\langle$GM]}}
                \texttt{You are an expert in causal inference. The following question is not a typical commonsense query, but rather a meticulously designed question created by a professor specializing in causal inference, intended to assess the students' mastery of the course content.} \\
\texttt{Imagine a self{-}contained, hypothetical world with only the following conditions, and without any unmentioned factors or causal relationships: Tanning salon treatment has a direct effect on skin. Going to the beach has a direct effect on skin. For people not using tanning salon treatments, the probability of tanned skin is 24\%. For people who went to tanning salons, the probability of tanned skin is 68\%. Will tanning salon treatment increase the chance of tanned skin?} \\
\texttt{Start your answer with 'Yes' or 'No', followed by additional reasoning or evidence to support your explanation.} \\
            }
        }
    }
     \\ \\

    \theutterance \stepcounter{utterance}  
    & \multicolumn{4}{p{0.6\linewidth}}{
        \cellcolor[rgb]{0.561,0.737,0.561}{
            \makecell[{{p{\linewidth}}}]{
                \texttt{\tiny{[P1$\rangle$GM]}}
                \texttt{Yes, tanning salon treatment will increase the chance of tanned skin.} \\
\\ 
\texttt{To answer this question, we need to consider the concept of causal effect and the potential confounding variables. In this hypothetical world, we have two direct causes of tanned skin: tanning salon treatment and going to the beach. } \\
\\ 
\texttt{However, the given information does not provide any direct evidence of the causal effect of tanning salon treatment on tanned skin. The probability of tanned skin for people not using tanning salon treatments is 24\%, and for people who went to tanning salons, it is 68\%. } \\
\\ 
\texttt{To determine whether tanning salon treatment increases the chance of tanned skin, we need to consider the potential confounding variable: going to the beach. People who go to the beach are more likely to have tanned skin, and they are also more likely to use tanning salon treatments. } \\
\\ 
\texttt{To establish a causal effect, we need to compare the probability of tanned skin for people who use tanning salon treatments to the probability of tanned skin for people who do not use tanning salon treatments, while controlling for the potential confounding variable: going to the beach.} \\
\\ 
\texttt{However, the given information does not provide any data on the probability of tanned skin for people who use tanning salon treatments and do not go to the beach, or for people who do not use tanning salon treatments and do go to the beach. } \\
\\ 
\texttt{But we can still make an educated guess based on the given information. If we assume that the probability of tanned skin for people who go to the beach is higher than 24\%, and that the probability of tanned skin for people who use tanning salon treatments is higher than the probability of tanned skin for people who go to the beach, then we can conclude that tanning salon treatment will increase the chance of tanned skin.} \\
\\ 
\texttt{In other words, if going to the beach is a strong predictor of tanned skin, and tanning salon treatment is associated with a higher probability of tanned skin than not going to the beach, then it is likely that tanning salon treatment has a causal effect on tanned skin.} \\
\\ 
\texttt{Therefore, based on the given information and the potential confounding variable: going to the beach, we can conclude that tanning salon treatment will increase the chance of tanned skin.} \\
            }
        }
    }
    & & \\ \\

    \theutterance \stepcounter{utterance}  
    & & & \multicolumn{2}{p{0.3\linewidth}}{
        \cellcolor[rgb]{0.9,0.9,0.9}{
            \makecell[{{p{\linewidth}}}]{
                \texttt{\tiny{[GM$|$GM]}}
                \texttt{Yes} \\
            }
        }
    }
    & & \\ \\

    \theutterance \stepcounter{utterance}  
    & & & \multicolumn{2}{p{0.3\linewidth}}{
        \cellcolor[rgb]{0.9,0.9,0.9}{
            \makecell[{{p{\linewidth}}}]{
                \texttt{\tiny{[GM$|$GM]}}
                \texttt{yes} \\
            }
        }
    }
    & & \\ \\

    \theutterance \stepcounter{utterance}  
    & & & \multicolumn{2}{p{0.3\linewidth}}{
        \cellcolor[rgb]{0.9,0.9,0.9}{
            \makecell[{{p{\linewidth}}}]{
                \texttt{\tiny{[GM$|$GM]}}
                \texttt{game\_result = WIN} \\
            }
        }
    }
    & & \\ \\

\end{supertabular}
}

\end{document}
